\section{研究背景与问题定义}

\subsection{研究背景}

建筑行业能耗占全球终端能源消费约30\%，碳排放占总量约28\%\cite{unep2020global}。中国在"双碳"目标（2030年碳达峰、2060年碳中和）背景下，既有建筑存量巨大且能耗水平参差不齐，如何精准识别高耗能建筑并制定科学的节能改造初步筛选清单，是城市管理者面临的核心挑战。

本研究依托西雅图市建筑能耗对标数据集（Seattle Building Energy Benchmarking, 2015--2024），该数据以``建筑--年份''为观测单元共计34,699条记录，含47个字段，覆盖能耗强度（SiteEUI）、温室气体排放（GHG）、建筑物理属性（建筑面积、层数、建造年份）、能源之星评分等信息。同时引入西雅图市按建筑类型和年份汇集的基准性能范围数据（Benchmarking Performance Ranges by Building Type），为能耗对标评估提供行业参照。

\subsection{问题定义}

城市建筑节能改造初步筛选面临三个核心问题：
\begin{enumerate}
    \item \textbf{能耗基准估计问题：}仅基于建筑静态属性（类型、面积、建造年代、区域等）和年度背景信息等预测时可获得的数据，能否合理估计建筑在特定年度的SiteEUI基准水平？
    \item \textbf{异常高耗能建筑识别问题：}哪些建筑实际能耗显著高于同类建筑的合理预期（模型预测基准和同类型统计基准）？这些建筑是否在近三年持续高耗能？
    \item \textbf{优先级排序问题：}在缺乏详细改造成本数据的条件下，如何综合预测残差、同类基准超耗、持续高耗能表现、碳排放强度、建筑老化程度等多维因素，构建建筑级节能改造优先级评分，并通过稳定性分析验证排序的可靠性？
\end{enumerate}

\subsection{数据说明}

本研究所用数据为公开的西雅图建筑能耗对标数据集，以"建筑--年份"为观测单元，共34,699条建筑年度记录，唯一建筑数量需逐年统计。同时使用西雅图市发布的按建筑类型与年份划分的基准性能范围数据，为计算相对基准的超耗比例提供参照。

